\documentclass[10pt,a4paper]{amsart}
\usepackage[margin=19mm]{geometry}
\usepackage{amsmath,amssymb,mathtools}
\usepackage[scr=rsfs,cal=euler]{mathalfa}
\usepackage{xfrac}
\usepackage[unicode,hidelinks,bookmarks=false]{hyperref}
\newcommand{\ie}{i.e.\ }
\newcommand{\cf}{cf.\ }
\newcommand{\resp}{respectively\ }
\newcommand{\Subsection}{\subsection}
\providecommand{\nobreakdash}{\nobreak\hbox{-}}
\setlength{\emergencystretch}{2em}
\multlinegap0pt
\usepackage[shortlabels]{enumitem}
\usepackage{amssymb}
\usepackage{csquotes}
\newtheorem{lemma}{Lemma}
\newtheorem{claim}{Claim}
\def\forces{\Vdash}
\newcommand{\la}{\langle}
    \newcommand{\one}{\mathop{1\hskip-3pt {\rm l}}}
\newcommand{\po}{\mathbb{P}}
\newcommand{\qo}{\mathbb{Q}}
\newcommand{\ra}{\rangle}
\renewcommand{\restriction}{\mathbin\upharpoonright}    % by default in amssymb it's mathrel
\def\s{\subseteq}
\DeclareMathOperator{\supp}{supp}
\title{The number of measures on very large measurable cardinals}
\author{Arthur W. Apter, Eyal Kaplan, Alejandro Poveda}
\date{Source: arXiv:2603.11710v1 (CC BY 4.0)}
\begin{document}
\maketitle

\noindent\emph{Source and licence.} The number of measures on very large measurable cardinals. Version arXiv:2603.11710v1 (CC BY 4.0), distributed by arXiv under the Creative Commons Attribution 4.0 International licence, http://creativecommons.org/licenses/by/4.0/, as recorded in the arXiv licence metadata for this exact version and shown on the abstract page https://arxiv.org/abs/2603.11710v1. Full licence notice: \texttt{licenses/arxiv-2603.11710-CC-BY-4.0.txt}.\par\smallskip

\noindent\textit{Editorial scope.} Counted unit: the article's lemma lemma (source0.tex lines 370--381). The article's complete original proof of it is delivered with it (source0.tex lines 383--464, one \texttt{proof} environment of 82 lines). No mathematics is written, repaired or re-derived: the statement and the proof are the article's own lines, in the article's own notation, and no retained line is substituted or re-flowed.  No further source-local context is needed: every cross-reference of every retained line resolves inside the two delivered spans.  Nothing was omitted from any retained span, so the package carries no omission record.  No retained line contains a citation, so the wrapper carries no bibliography and no bibliography entry.  No retained line contains a figure environment, an image-inclusion command or drawing code, so no image is delivered and the package declares no asset dependency.  Editorial formatting only, in addition: an AMS \texttt{amsart} preamble that restates the article's own declarations and macros used by the retained lines, namely the environments \texttt{lemma}, \texttt{claim} on separate counters, together with the standard packages those lines need.  The retained environments print in this document as Lemma 1, Claim 1 in order of appearance, whereas the article prints the same items with the numbering of its own full document.  The complete native source of the article as distributed in the arXiv e-print for this exact version is bundled unchanged as \texttt{sources/196232/source0.tex}.
\begin{lemma}[The fusion lemma]\label{Lemma: fusion lemma}
    Assume that $\kappa$ is Mahlo and $I\subseteq \kappa$ is a stationary set of inaccessible cardinals. Let $\la \po_\alpha, \dot{\qo}_\alpha \colon \alpha<\kappa \ra$ be an $I$-spaced nonstationary support iterated forcing, such that:
    \begin{enumerate}
        \item For every $\alpha<\kappa$, $\mathrm{rank}(\dot{\qo}_\alpha) < \min(I \setminus (\alpha+1))$.
        \item For every $\alpha< \kappa$, $\one \Vdash_{\po_\alpha} {``\dot{\qo}_\alpha \text{ is }\alpha\text{-strategically closed}}$".
        \item For every $\alpha\in \kappa\setminus I$, $\one \Vdash_{\po_\alpha} {``\dot{\qo}_\alpha \text{ is trivial forcing".}}$
       
    \end{enumerate}
    Let $\po$ be the nonstationary support limit of $\la \po_\alpha \colon \alpha<\kappa \ra$. Then for every $p\in \po$ and a sequence $\vec{d} = \la d(\alpha)\colon \alpha<\kappa \ra$ of dense open subsets of $\po$, there exists a condition $p^*\leq p$ and a club $C\subseteq \kappa$ such that for every $\alpha\in C$, 
    $$\{ r\in \po_{\alpha+1} \colon r^{\smallfrown} (p^*\setminus \alpha+1) \in d(\alpha) \}$$
    is dense in $\po_{\alpha+1}$ below $p^*\restriction (\alpha+1)$.
\end{lemma}

\begin{proof}
    Fix for every $\alpha\in I$ a $\po_\alpha$-name $\dot{\tau}_\alpha$ for a strategy of Player II in the game $G_{\alpha}( \dot{\qo}_\alpha)$.
    We construct sequences:
    \begin{itemize}
        \item $\la  p_i \colon i<\kappa\ra$ a decreasing  sequence of conditions in $\po$.
        \item $\la \alpha_i \colon i<\kappa \ra$ a continuous, increasing cofinal sequence in $\kappa$.
        \item $\la C_i \colon i<\kappa \ra$ a decreasing sequence of club subsets of $\kappa$ (with respect to inclusion), where each $C_i$ is disjoint from $\supp(p_i)$. 
    \end{itemize}
    The construction is done in such a way that the following properties hold:
    \begin{itemize}
        \item For every $i<j<\kappa$, $p_{i}\restriction \alpha_i+1 = p_j \restriction \alpha_i+1$.
        \item For every $i<\kappa$, $\{ r\in \po_{\alpha_{i}+1} \colon r^{\frown} p_{i}\setminus (\alpha_i+1) \in d(\alpha_i) \}$ is dense in $\po_{\alpha_i+1}$ below $p_i\restriction (\alpha_i+1)$.
        \item For every $i<j$, $\alpha_j \in C_i$. 
        \item For every $i<\kappa$ and $\alpha\in \supp(p_i)\setminus (\alpha_i+1)$, $p_i\restriction \alpha$ forces that $\la p_{j}(\alpha) \colon j<i, \alpha\in \supp(p_j) \ra$ is the sequence of moves of Player II in a partial run of the game $G_\alpha(\dot{\qo}_\alpha)$ in which Player II plays according to their winning strategy $\dot{\tau}_\alpha$.
    \end{itemize}
    
    Start the construction by letting $p_0 = p$, $\alpha_0 = 0$, and $C_0$ be any club disjoint from $\supp(p)$. Assume that $p_i, \alpha_i, C_i$ have been constructed for some $i<\kappa$. Let $\alpha_{i+1} = \min(C_i \setminus \alpha_i+1)$. We define $p_{i+1}\leq p_i$ such that:
    \begin{enumerate}
        \item $p_{i+1}\restriction \alpha_{i+1} = p_i \restriction \alpha_{i+1}$.
        \item $\alpha_{i+1}\notin \supp(p_{i+1})$ (and since $\alpha_{i+1}\in C_i$ and $C_i$ is disjoint from $\supp(p_i)$,  $p_{i+1}\restriction (\alpha_{i+1}+1) = p_{i}\restriction (\alpha_{i+1}+1)$).
        \item $p_{i+1}\setminus (\alpha_{i+1}+1)$ satisfies the following two properties:
        \begin{enumerate}
            \item $p_{i+1}\setminus (\alpha_{i+1}+1)$ is forced by $p_{i+1}\restriction (\alpha_{i+1}+1)$ to be an  extension $s\leq p_i\setminus (\alpha_{i+1}+1)$ for which $\{ r\in \po_{\alpha_{i+1}+1} \colon  r^{\frown} s\in d(\alpha_{i+1}) \}$ is a dense subset of $\po_{\alpha_{i+1}+1}$ below $p\restriction (\alpha_{i+1}+1)$.
            \item For every $\alpha\in \supp(p_{i+1})\setminus (\alpha_{i+1}+1)$, $p_{i+1}\restriction \alpha$ forces that $\la p_j(\alpha) \colon j\leq i, \alpha\in \supp(p_j) \ra$ is the sequence of moves of Player II in a run of the game $G_\alpha(\dot{\qo}_\alpha)$ in which Player II plays according to $\dot{\tau}_\alpha$.
        \end{enumerate}
    \end{enumerate}
\begin{claim}\label{Claim: fusion lemma, arranging clause 3}
    Clause~(3) can be arranged.
\end{claim}
\begin{proof}[Proof of claim]
   We need to argue that such a condition $s$ exists.  This essentially follows from the fact that $|\po_{\alpha_{i+1}+1}|$ is strictly below the amount of strategic closure of $\po\setminus (\alpha_{i+1}+1)$.\footnote{Indeed, the argument we provide below shows that for every $\alpha<\kappa$, the forcing $\po\setminus (\alpha+1)$ is $ \min(I\setminus (\alpha+1))$-strategically closed.} More formally, let $\langle q_\alpha\colon \alpha<\chi\rangle$ be an injective enumeration of all conditions in $\mathbb{P}_{\alpha_{i+1}+1}$ below $p_i\restriction (\alpha_{i+1}+1)$. We construct $\mathbb{P}_{\alpha_{i+1}+1}$-names $\langle \dot{p}_{i+1}^\beta\colon \beta\leq \chi\rangle$ as follows:
   \begin{itemize}
       \item $p_{i+1}\restriction (\alpha_{i+1}+1)$ forces that $\la \dot{p}^\beta_{i+1} \colon \beta\leq\chi \ra$ is a decreasing sequence of conditions in $\mathbb{P}\setminus (\alpha_{i+1}+1).$
       \item For every $\beta<\chi$, ${q_\beta}^{\frown} \dot{p}^{\beta}_{i+1}\in d(\alpha_{i+1})$.
       \item For every $\beta<\chi$ and $\alpha\in \supp(\dot{p}^\beta_{i+1})\setminus (\alpha_{i+1}+1)$, $\dot{p}^{\beta}_{i+1}\restriction \alpha$ forces that $\la \dot{p}^\gamma_{i+1}(\alpha) \colon \gamma<\beta, \alpha\in \supp(\dot{p}^\gamma_{i+1}) \ra$ is the sequence of moves of Player II in a run of the game $G_\alpha(\dot{\qo}_\alpha)$ in which player II plays according to the strategy $\dot{\tau}_\alpha$. 
   \end{itemize}
   Once such a sequence is constructed, we pick the desired $\po_{\alpha_{i+1}+1}$-name $\dot{s}$ to be a condition on $\po\setminus \alpha_{i+1}+1$ with $\supp(\dot{s}) = \bigcup_{\beta\leq \chi} \supp(\dot{p}^\beta_{i+1})$, as follows: Assume that $\alpha\in (\alpha_{i+1}, \kappa)$ and $\dot{s}\restriction \alpha$ has been constructed. Since $\dot{s}\restriction \alpha$ extends ${p}_{i}\restriction (\alpha_{i+1}+1,\alpha)$, it forces that $\la p_{j}(\alpha) \colon j\leq i, \alpha\in \supp(p_j)\ra$ is the sequence of moves of Player II in the game $G_\alpha(\dot{\qo}_\alpha)$ according to $\dot{\tau}_\alpha$. Assume that, in the same game, Player I picks the condition $\dot{p}^\chi_{i+1}(\alpha)$ at stage $i+1$ (and, if $\alpha\notin \supp(\dot{p}^\chi_{i+1})$, we assume that Player I picks $\one_{\dot{\qo}_\alpha}$). Let $\dot{s}(\alpha)$ be the condition forced by $\dot{s}\restriction \alpha$ to be the response of Player II when playing according to $\dot{\tau}_\alpha$. Since $\alpha$ is an inaccessible cardinal strictly above $\alpha_{i+1}$ (and in particular,  $\alpha>\chi$), the strategy $\dot{\tau}_\alpha$ ensures that Player II has a valid response at each stage in the run of the game described above (recall that $\dot{\tau}_\alpha$ witnesses that $\dot{\qo}_\alpha$ is $\alpha$-strategically closed). This concludes the definition of the condition $\dot{s}$. Since $\chi < \min(I\setminus (\alpha_{i+1}+1))$, the set $\supp(\dot{s})$ is forced to be nonstationary in inaccessibles of $I\setminus (\alpha_{i+1}+1)$, so $\dot{s}\in \po\setminus (\alpha_{i+1}+1)$. It is  not hard to verify that it satisfies the requirements in clause $(3)$ above. 
   
   
   Thus, it remains to construct the sequence $\langle \dot{p}_{i+1}^\beta\colon \beta\leq \chi\rangle$.

   We first describe how the condition $\dot{p}^0_{i+1}$ is chosen. For that, pick a condition $q'\leq {q_0}^{\frown} (p_i\setminus (\alpha_{i+1}+1))$ such that $q'\in d(\alpha_{i+1})$. Define $\dot{p}^0_{i+1}$ such that $\supp(\dot{p}^0_{i+1}) = \supp(q')$, and, for every $\alpha\in \supp(\dot{p}^0_{i+1})$, $\dot{p}^0_{i+1}\restriction \alpha$ forces that $\dot{p}^0_{i+1}(\alpha)$ is a name for the response dictated to Player II by $\dot{\tau}_\alpha$ after Player I picked $q'(\alpha)$ on their first move in the game $G_\alpha(\dot{\qo}_\alpha)$. 

   Assume now that $\beta\leq \chi$ and $\la \dot{p}^\gamma_{i+1} \colon \gamma<\beta \ra$ have been defined. Let $\dot{t}\in \po\setminus (\alpha_{i+1}+1)$ be inductively constructed, as follows: For every $\alpha\in (\alpha_{i+1}, \kappa)$, if $\dot{t}\restriction \alpha$ was defined, it forces that $\dot{t}(\alpha)$ is the move dictated to Player II by the strategy $\dot{\tau}_\alpha$ in a run of the game $G_{\alpha}(\dot{\qo}_\alpha)$ in which the moves of Player II were $\la \dot{p}^\gamma_{i+1}(\alpha) \colon \gamma<\beta, \alpha\in \supp(\dot{p}^\gamma_{i+1}) \ra$. Note that the condition $\dot{t}\in \po\setminus (\alpha_{i+1}+1)$ constructed this way is forced by $p_{i+1}\restriction (\alpha_{i+1}+1)$ to be a lower bound of $\la  \dot{p}^\gamma_{i+1} \colon \gamma<\beta\ra$. Pick an extension $q'\leq (q_{\gamma})^{\frown} \dot{t}$ such that $q'\in d(\alpha_{i+1})$. Now let $\dot{p}^{\beta}_{i+1}\in \po\setminus (\alpha_{i+1}+1)$ be such that, for every $\alpha\in (\alpha_{i+1}+1, \kappa)$, $\dot{p}^\beta_{i+1}\restriction \alpha$ forces that $\dot{p}^\beta_{i+1}(\alpha)$ is the move dictated to Player II by $\dot{\tau}_\alpha$ in the following run of the game $G_\alpha(\dot{\qo}_\alpha)$: The game begins with the run in which the moves of Player II were $\la \dot{p}^\gamma_{i+1}(\alpha) \colon \gamma<\beta, \alpha\in \supp(\dot{p}^\gamma_{i+1}) \ra^{\frown} \dot{t}(\alpha)$. Then Player I replies with the condition $q'(\alpha)$, and the response of Player II in its next round according to $\dot{\tau}_\alpha$ is taken to be $\dot{p}^\beta_{i+1}(\alpha)$.  

   This concludes the inductive construction of  the sequence $\la \dot{p}^\beta_{i+1} \colon \beta\leq \chi \ra$, and therefore concludes also the construction of the condition $\dot{s}$.
 
\end{proof}
    Let us argue that the condition $p_{i+1}$ constructed this way indeed has a nonstationary support. Fix an inaccessible $\lambda>\alpha_{i+1}$. We claim that there exists a club in $\lambda$ that belongs to $V$, and is forced by $p_{i+1}\restriction (\alpha_{i+1}+1)$ to be disjoint from $\supp(p_{i+1}\setminus (\alpha_{i+1}+1))$. Indeed, there exists a $\po\restriction (\alpha_{i+1}+1)$-name $\dot{D}$ for a club in $\lambda$ which is disjoint from $\supp(p_{i+1}\setminus (\alpha_{i+1}+1)$. By the fact that $\po\restriction (\alpha_{i+1}+1)$ is small relative to $\lambda$ (and, in particular, is $\lambda$-c.c.), there exists a club $D^*\in V$ such that $p_{i+1}\restriction (\alpha_{i+1}+1)$ forces that $D^*$ is contained in $\dot{D}$; in particular, $D^*$ is disjoint from $\supp(p_{i+1})$.
    
    Finally, let $C_{i+1}\subseteq C_i$ be a club subset of $\kappa$ in $V$ which is  disjoint from the support of $p_{i+1}$. This concludes the successor step in the construction.

    For limit steps in the construction, assume that $\la p_j, \alpha_j , C_j \colon j<i \ra$ were constructed for some $i<\kappa$. Let $\alpha_i = \sup_{j<i}\alpha_j$. We define $p_i \in \po$ such that:
    \begin{enumerate}
     \item $p_i\restriction \alpha_{i} = \bigcup_{j<i} p_j \restriction (\alpha_j+1)$. Note that the union is increasing, and in the case where $\alpha_i$ is inaccessible, the sequence $\la \alpha_j \colon j<i \ra$ is a club in $\alpha_i$ disjoint from the $\supp(p_i\restriction \alpha_i)$. Thus,  $p_i\restriction \alpha_i\in \po_{\alpha_i}$.
        \item $\alpha_i \notin \supp(p_i)$. In particular, $p_i\restriction (\alpha_i+1 )\leq p_j \restriction (\alpha_i+1)$ for every $j<i$, since $\alpha_i$ lies outside $\supp(p_j)$  in that of a  limit point of $C_j$.
        \item $p_{i}\setminus (\alpha_{i}+1)$ satisfies the following two properties:
        \begin{enumerate}
            \item $p_i\setminus (\alpha_i+1)$ is forced by $p_i\restriction (\alpha_i+1)$ to be a common extension $\dot{s}$ of $\la p_j \setminus (\alpha_i+1) \colon j<i \ra$ for which $\{  r\in \po_{\alpha_i+1} \colon r^{\frown} \dot{s} \in d(\alpha_i) \}$ is a dense subset of $\po_{\alpha_i+1}$ below $p\restriction (\alpha_{i}+1)$.
            \item For every $\alpha\in \supp(p_{i}\setminus (\alpha_{i}+1))$, $p_{i}\restriction \alpha$ forces that $$\la p_j(\alpha) \colon j\leq i, \alpha\in \supp(p_j) \ra$$ is the sequence of moves of Player II in a run of the game $G_\alpha(\dot{\qo}_\alpha)$ in which Player II plays according to $\dot{\tau}_\alpha$.
        \end{enumerate}
        The proof that such $\dot{s}$ can be constructed is similar to the proof of Claim \ref{Claim: fusion lemma, arranging clause 3}. Note that for each coordinate $\alpha > \alpha_i$, the forcing $\dot{\qo}_\alpha$ is sufficiently strategically closed in order to pick local lower bounds.
    \end{enumerate}

    
    Finally, let $p^* =\bigcup_{i<\kappa} p_i\restriction (\alpha_i+1)$. {Clearly, $\supp(p^*)\cap \lambda$ is nonstationary in $\lambda$, for every inaccessible cardinal $\lambda<\kappa$.} Also, $\supp(p^*)$ is nonstationary in $\kappa$, as witnessed by the club  $C =\langle \alpha_i \colon i<\kappa \rangle$. To complete the proof of Lemma \ref{Lemma: fusion lemma}, we have to verify the following:
    \begin{claim}
     For every $\alpha\in C$, 
    $$\{ r\in \po_{\alpha+1} \colon r^{\smallfrown} (p^*\setminus \alpha+1) \in d(\alpha) \}$$
    is dense in $\po_{\alpha+1}$ below $p^*\restriction (\alpha+1)$.
    \end{claim}
    \begin{proof}[Proof of claim]
        Let $\alpha\in C$. Our inductive construction ensures that $$D=\{r\in \mathbb{P}_{\alpha+1}\colon r\leq p_\alpha\restriction (\alpha+1)\,\wedge\, r{}^\smallfrown p_{\alpha}\setminus (\alpha+1)\in d(\alpha)\}$$
        is dense below $p_{\alpha}\restriction (\alpha+1)=p^*\restriction(\alpha+1).$
        
        Let $r\in D$ be arbitrary. Since $p^*\leq p_\alpha$, in particular,  $$p^*\restriction (\alpha+1)\forces p^*\setminus (\alpha+1)\leq p_\alpha\setminus (\alpha+1).$$
        But then $r{}^\smallfrown p^*\setminus (\alpha+1)\leq r{}^\smallfrown p_\alpha\setminus (\alpha+1)$, which yields
        $$D\s \{r\in \mathbb{P}_{\alpha+1}\colon r\leq p^*\restriction(\alpha+1)\,\wedge\, r{}^\smallfrown p^*\setminus (\alpha+1)\in d(\alpha)\},$$
        and as a result the latter is dense below $p^*\restriction (\alpha+1).$
    \end{proof}
    This concludes the proof of Lemma \ref{Lemma: fusion lemma}.
\end{proof}

\end{document}
